\documentclass{article}
\usepackage[margin=2cm]{geometry}
\usepackage{amsmath}
\usepackage{graphicx}
\usepackage{booktabs}
\usepackage[utf8]{inputenc}
\usepackage{ragged2e}
\setlength{\emergencystretch}{3em}
\begin{document}
\sloppy

\noindent\hfill\includegraphics[height=1.8em]{assets/figure_header_logo.png}\par\noindent\hrule\vspace{4pt}

in the limit the three punctured sphere we have defined has all Cx vanishing but the one corresponding to the constant polynomial. The Koornwinder
polynomialshavehigher rank generalizationswhich should be relevantfor
higher rank E string theories.In those cases we do notknow the three punctured spheres and the relation to Koornwinder polynomials can provide a
useful tool to study the indices of these models.The limit we considered
does not have a special physical meaning a priori, however the fact that
the expressions become simple and the fact that one might generalize the
discussion to the higher rank case, make the limit of potential interest.


\subsection*{A Index definitions}


We compute the supersymmetric index using the standard definitions of .
The index of chiral field charged under favor U(1) symmetry with charge S
and having R-charge  is


\begin{equation}
\Gamma _ { e } \big ( ( q p ) ^ { \frac { \mathfrak { R } } { 2 } } u ^ { S } \big ).
\end{equation}


The parameter u is fugacity for the favor symmetry. We define here


\begin{equation}
\Gamma _ { e } ( u ) = \prod _ { i, j = 0 } ^ { \infin } \frac { 1 - \frac { 1 } { u } q ^ { i + 1 } p ^ { j + 1 } } { 1 - u p ^ { i } q ^ { j } }.
\end{equation}


We will use the following definitions


\begin{equation}
( s ; q ) = \prod _ { i = 1 } ^ { \infin } \big ( 1 - s q ^ { i - 1 } \big ), \qquad \theta _ { r } ( u ) = \prod _ { j = 1 } ^ { \infin } \big ( 1 - u r ^ { j - 1 } \big ) \big ( 1 - r ^ { j } / u \big ).
\end{equation}


Finally we use the condensed conventions


Contour integrals in the paper are around the unit circle unless we state
otherwise.


\subsection*{B Computation of the sphere with two punctures}


We give here the derivation of equation. The computation involves calculatingseveralcontourintegralsoverproductsofellipticgammafunctions

\end{document}
